\chapter{问题描述与实验设计}\label{chap:introduction}

\section{八皇后问题形式化定义}

设棋盘规模为 $N\times N$（本实验取 $N=8$）。定义状态向量
\begin{equation}
  \mathbf{q}=(q_0,q_1,\ldots,q_{N-1}),\quad q_i\in\{0,1,\ldots,N-1\},
\end{equation}
其中 $q_i$ 表示第 $i$ 列皇后所在行。合法解需满足：
\begin{align}
  q_i       & \neq q_j,\quad   &  & \forall\ 0\le i<j<N, \\
  |q_i-q_j| & \neq |i-j|,\quad &  & \forall\ 0\le i<j<N.
\end{align}
前者保证不同行，后者保证不在同一对角线。

\section{方法对照设计}

报告将方法组织为六个可比较的实现，并与任务要求建立一一映射，如\tabref{tab:method-overview}所示。

\begin{table}[htbp]
  \centering
  \caption{求解方法与任务书条目对应关系}\label{tab:method-overview}
  \begin{tabular}{p{2.0cm}p{4.4cm}p{6.4cm}}
    \toprule
    方法编号 & 方法名称                    & 对应课程实验要求          \\
    \midrule
    M1   & 命令式穷举+逻辑筛选              & 先非逻辑穷举，再逻辑筛选      \\
    M2   & 逻辑求解（\texttt{permuteq}） & 不使用命令式穷举，仅逻辑描述    \\
    M3   & 逻辑求解（自定义不等关系）           & 补充库中缺失函数并参与求解     \\
    M4   & 深度优先搜索（DFS）             & 使用图搜索求解并比较        \\
    M5   & 广度优先搜索（BFS）             & 使用图搜索求解并比较        \\
    M6   & 位运算回溯                   & 课程扩展实现，用于对比工程优化效果 \\
    \bottomrule
  \end{tabular}
\end{table}

\section{实验环境与复现流程}

实验环境配置如下：
\begin{itemize}
  \item 操作系统：Windows 11；
  \item Python 版本：3.12.7；
  \item 逻辑编程库：Kanren 0.2.3；
  \item 实验脚本：\texttt{experiments/eight\_queens\_benchmark.py}；
  \item 结果文件：\texttt{experiments/eight\_queens\_benchmark.json}。
\end{itemize}

接口说明与基础语法主要参考 Python 官方文档、Kanren 项目文档及课程提供资料\cite{python_docs,logpy_repo,runoob_python}。

复现实验采用以下固定流程：
\begin{enumerate}
  \item 设定参数 $N=8$，每种方法独立运行；
  \item 每种方法重复 3 次，记录均值、标准差、最小值、最大值；
  \item 比较六种方法的解数量是否一致（应为 92）；
  \item 记录状态展开规模并结合复杂度进行解释。
\end{enumerate}

\section{评价指标}

评价指标分为两类：
\begin{itemize}
  \item \textbf{正确性指标}：解数量、最小字典序解、最大字典序解是否一致；
  \item \textbf{性能指标}：平均耗时与展开状态数。
\end{itemize}

其中耗时通过 \texttt{time.perf\_counter()} 统计，单位为毫秒。

\section{复杂度分析基线}

从理论上看，若不利用剪枝，八皇后搜索空间会快速膨胀。作为后续对照基线，可给出如下上界描述：
\begin{itemize}
  \item M1、M2：基于全排列候选，时间上界可记为 $O(N!\cdot N^2)$；
  \item M3：仍属于指数级搜索，但通过增量约束可显著减少无效分支；
  \item M4、M5：最坏情况下均为指数级，DFS 空间压力通常小于 BFS\cite{cormen2022}；
  \item M6：搜索树与回溯法同阶，但位运算降低了单节点常数开销。
\end{itemize}

第\ref{chap:math}章将用实测数据验证这些判断。
